% Generated draft: requires mathematical and source review.
\documentclass[11pt]{article}
\usepackage[margin=25mm]{geometry}
\usepackage{fontspec}
\setmainfont{Latin Modern Roman}
\usepackage{amsmath,amssymb,mathtools,unicode-math}
\setmathfont{Latin Modern Math}
\usepackage{enumitem,needspace,xurl,hyperref}
\setlength{\parindent}{0pt}
\setlength{\parskip}{5pt}
\setlength{\emergencystretch}{3em}
\newcommand{\N}{\mathbb{N}}
\newcommand{\Z}{\mathbb{Z}}
\newcommand{\Q}{\mathbb{Q}}
\newcommand{\R}{\mathbb{R}}
\newcommand{\C}{\mathbb{C}}
\newcommand{\F}{\mathbb{F}}
\newcommand{\A}{\mathbb{A}}
\newcommand{\PP}{\mathbb{P}}
\newcommand{\E}{\mathbb{E}}
\newcommand{\eps}{\varepsilon}
\newcommand{\dd}{\,\mathrm{d}}
\begin{document}
\section*{UnsolvedMath ID: 30001001}
\textbf{Source problem code:} catalogue entry
\subsection{Definitions and mathematical statement}
Fix the matrix norm intended by the source, and let \[\mathcal V=\{A\in M_2(\R):\det A=2\},\qquad B_T=\{A\in\mathcal V:\|A\|\leq T\}.\] Let $\mathcal P$ be the positive rational primes and define \[N_{\mathrm{prime}}(T)=\#\left\{\begin{pmatrix}a&b\\c&d\end{pmatrix}:a,b,c,d\in\mathcal P,\ ad-bc=2,\ \left\|\begin{pmatrix}a&b\\c&d\end{pmatrix}\right\|\leq T\right\}.\] Entries need not be distinct. The matrices are counted as ordered integral matrices, without quotienting by row/column operations or group actions. The qualitative question is whether infinitely many such matrices exist, equivalently $N_{\mathrm{prime}}(T)\to\infty$ as $T\to\infty$. The wording infinitely many in $B_T$ cannot literally apply to one fixed finite $T$, since a bounded norm ball contains only finitely many integral matrices; it is interpreted in the growing-ball sense explicitly indicated by the quantitative question. The stronger proposed prime matrix theorem is \[\exists C>0:\qquad N_{\mathrm{prime}}(T)\sim C\frac{T^2}{(\log T)^4}\qquad(T\to\infty),\] meaning the ratio tends to one. The norm is not specified, and $C$ is allowed to depend on that fixed norm; no numerical or local-density formula for it is supplied. The determinant is exactly positive two, not its absolute value or a varying determinant. The supplied 2026-08-21 triage labels the question open, says the report explicitly calls it open, and reports no later resolution found.
\subsection{Short English statement}
Are there infinitely many determinant-two integral matrices with four positive prime entries, and does their count in a growing norm ball have the stated $T^2/(\log T)^4$ asymptotic? The source's fixed-ball wording is read explicitly as a limit through growing balls.
\subsection{Sources}
\begin{itemize}\item Local UnsolvedMath registry, permanent ID 30001001, source problem catalogue entry, \emph{Prime-Entry Matrices of Fixed Determinant}, proposed year 2008; no individual proposer supplied.\item Supplied citation: \emph{Real Analysis, Harmonic Analysis and Applications} (2009), Oberwolfach Reports, DOI 10.4171/owr/2008/32, pp.\ 1771--1850: \url{https://doi.org/10.4171/owr/2008/32}. Supplied literature label: original Oberwolfach report.\end{itemize}
\end{document}
